% ============================================================
% CV — Pierre PUGET (sans photo)
% Compilation : pdflatex CV_Puget_Pierre.tex
% Packages standards uniquement — aucun fichier externe requis.
% ============================================================
\documentclass[11pt, a4paper]{article}

\usepackage[a4paper, top=1.6cm, bottom=1.6cm, left=1.9cm, right=1.9cm]{geometry}
\usepackage[utf8]{inputenc}
\usepackage[T1]{fontenc}
\usepackage[french]{babel}
\usepackage{xcolor}
\usepackage{titlesec}
\usepackage{enumitem}
\usepackage{hyperref}
\usepackage{ragged2e}
\usepackage{array}

\pagestyle{empty}
\setlength{\parindent}{0pt}

% --- Palette (sobre : encre + un seul accent) ---
\definecolor{ink}{HTML}{1A1A1A}
\definecolor{muted}{HTML}{5B5B5B}
\definecolor{accent}{HTML}{2E5266}
\definecolor{rule}{HTML}{D8D8D8}

\color{ink}
\hypersetup{colorlinks=true, urlcolor=accent, linkcolor=accent}

% --- Titres de section ---
\titleformat{\section}
    {\normalfont\large\bfseries\color{accent}}
    {}{0pt}{}
    [\vspace{2pt}{\color{rule}\titlerule[0.8pt]}]
\titlespacing*{\section}{0pt}{14pt}{8pt}

% --- Entrée générique (diplôme / expérience) ---
% #1 titre   #2 dates   #3 lieu   #4 description (peut être vide)
\newcommand{\cventry}[4]{%
    \noindent\begin{tabular*}{\linewidth}{@{}l@{\extracolsep{\fill}}r@{}}
        \textbf{#1} & {\small\color{muted}#2}\\
    \end{tabular*}\par
    {\small\color{accent}#3}\par
    \ifx&#4&%
    \else
        \vspace{2pt}{\small\color{ink}#4}\par
    \fi
    \vspace{7pt}
}

% --- Ligne de contact (icônes texte, aucune police externe) ---
\newcommand{\cvcontact}[2]{%
    \textbf{#1}~#2
}

\begin{document}

% ============================================================
% EN-TÊTE
% ============================================================
\begin{center}
    {\Huge\bfseries Pierre Puget}\\[4pt]
    {\large\color{muted}Étudiant en BUT Informatique --- Développement \& Systèmes}\\[8pt]
    {\small
        \cvcontact{Email~:}{\href{mailto:pugetpierre07@gmail.com}{pugetpierre07@gmail.com}}
        \quad\textcolor{rule}{|}\quad
        \cvcontact{Tél.~:}{06~68~35~09~17}
        \quad\textcolor{rule}{|}\quad
        \cvcontact{Adresse~:}{55 route de Crau, 13200 Arles}
        \quad\textcolor{rule}{|}\quad
        \cvcontact{Permis~:}{B}
    }
\end{center}
\vspace{4pt}
{\color{rule}\hrule height 0.8pt}
\vspace{10pt}

\RaggedRight
Actuellement en première année de BUT Informatique à l'IUT d'Arles, ce premier semestre a confirmé mon attrait pour les aspects techniques et la logique algorithmique. Passionné par ces enjeux, je souhaite orienter mon parcours vers la cybersécurité.

% ============================================================
% DEUX COLONNES : contenu principal (gauche) / profil complémentaire (droite)
% ============================================================
\begin{minipage}[t]{0.63\textwidth}

\section*{Compétences informatiques}
\textbf{Développement et Algorithmique}\par
{\small\color{muted}Apprentissage de la logique de programmation et des algorithmes via les langages Python et C++.\par}
\vspace{5pt}
\textbf{Web et Données}\par
{\small\color{muted}Bases de l'intégration web (HTML, CSS) et initiation à la gestion de bases de données (SQL).\par}
\vspace{5pt}
\textbf{Systèmes et Environnements}\par
{\small\color{muted}Maîtrise des bases de l'administration système (Linux/Bash, Windows). Compréhension de l'architecture des ordinateurs et des systèmes d'exploitation (OS).\par}

\section*{Diplômes et Formations}

\cventry{Baccalauréat}{De 2024 à 2025}{Lycée Louis Pasquet, Arles --- Mention Assez Bien}{}
\cventry{BAFA}{2025}{IFAC, Nîmes}{}
\cventry{PSC1}{2025}{La Croix Rouge, Arles}{}
\cventry{Diplôme National du Brevet}{De 2018 à 2022}{Collège Vincent Van Gogh, Arles --- Mention Très Bien}{}

\end{minipage}%
\hfill
\begin{minipage}[t]{0.33\textwidth}

\section*{Langues}
\textbf{Anglais}\par
{\small\color{muted}Niveau B2 (Intermédiaire avancé)}\par
\vspace{6pt}
\textbf{Espagnol}\par
{\small\color{muted}Les bases (A2)}\par

\section*{Atouts}
\begin{itemize}[leftmargin=12pt, itemsep=3pt, topsep=2pt, label={\small\color{accent}\textbullet}]
    \item \small Curiosité technique
    \item \small Esprit logique et analytique
    \item \small Rigueur
    \item \small Bon relationnel
    \item \small Sens de l'organisation
    \item \small Esprit d'équipe
\end{itemize}

\section*{Centres d'intérêt}
\begin{itemize}[leftmargin=12pt, itemsep=4pt, topsep=2pt, label={\small\color{accent}\textbullet}]
    \item \small Jeux vidéo
    \item \small Musique : guitare, piano
    \item \small Sport : ceinture noire de karaté, MMA, musculation, rugby
    \item \small Plongée : niveau 1
\end{itemize}

\end{minipage}

% ============================================================
% EXPÉRIENCES PROFESSIONNELLES (pleine largeur)
% ============================================================
\section*{Expériences professionnelles}

\cventry{Cours de soutien scolaire}{De 2024 à 2026}{À domicile}{Je donne des cours de soutien scolaire à deux élèves. Je les aide à faire leurs devoirs et je leur propose des jeux ludiques, ce qui leur a permis de beaucoup s'améliorer.}

\cventry{Animateur et Assistant Sanitaire}{Juillet 2025}{BAFA, Serre Chevalier}{Gestion d'un groupe de mineurs alliant conception d'activités ludiques et responsabilité du pôle sanitaire (soins, traitements, registre). Une mission menée avec succès grâce à une grande rigueur et un bon relationnel.}

\cventry{Stage de mathématiques}{2024}{Campus de Luminy, Marseille}{Participation à un stage encadré par des doctorants, axé sur la résolution de problèmes techniques en mathématiques, dans un contexte de découverte de la recherche scientifique.}

\cventry{Stage en entreprise}{2022}{ESISAR --- Grenoble INP, Valence}{Stage de troisième réalisé à l'ESISAR, école d'ingénieurs spécialisée, au cours duquel j'ai découvert les bases de l'informatique et de la cybersécurité à travers des échanges avec des étudiants et des professionnels du secteur.}

\end{document}
